%% kapa: exam=1 round=24 year=2013 period=2 form=A subject=appraisal_law subject_ko=부동산관계법규 seq=4 number=44
%% source: Q-Net 감정평가사 기출문제 (cst00309) artlSeq=5073802 file=extracted/1008/5073802/2013년 제24회 감정평가사 1차시험 기출문제/2교시(A형).hwp
%% conversion: kordoc 4.24.51 -> tools/gen1.py (auto); review=proofread-sonnet; flags=none
\begin{question}{44}
국토의 계획 및 이용에 관한 법령상 도시ㆍ군기본계획에 관한 설명으로 옳은 것은?
\begin{choices}
  \item 시장 또는 군수는 지역여건상 필요하다고 인정되면 인접한 시 또는 군의 관할구역의 전부를 포함하여 도시ㆍ군기본계획을 수립할 수 있다.
  \item 도시ㆍ군기본계획은 광역도시계획과 도시ㆍ군관리계획 수립의 지침이 되는 계획을 말한다.
  \item 시장 또는 군수가 도시ㆍ군기본계획을 변경하려면 국토교통부장관의 승인을 얻어야 한다.
  \item 광역도시계획이 수립되어 있는 지역의 도시ㆍ군기본계획의 내용이 그 광역도시계획의 내용과 다른 때에는 도시ㆍ군기본계획의 내용이 우선한다.
  \item 수도권에 속하는 인구 10만 이하의 시 또는 군의 경우는 도시ㆍ군기본계획을 수립하지 아니할 수 있다.
\end{choices}
\end{question}
\answer{1}
